\section{Exact approximation floor for the nonlinear stress task}
\label{app:nonlinear}
This analytic diagnostic was developed after the primary protocol was
frozen. It does not change training, tuning, or the original Monte Carlo
endpoint. Let $Z_1,Z_2,Z_3$ be independent standard normal coordinates after
an orthogonal rotation, and let
\[
 f_*(Z)=.8\tanh(1.5Z_1)\tanh(1.5Z_2)+.4\sin Z_3.
\]
Define
\[
 a=\mathbb E[Z\tanh(1.5Z)],\quad
 b=\mathbb E[\tanh^2(1.5Z)],\quad c=.8a^2.
\]
\begin{proposition}[A width-independent approximation obstruction]
The orthogonal projection of $f_*$ onto the class of all centered quadratic
forms plus constants is $cZ_1Z_2$. If $u_1,u_2$ are the corresponding
orthonormal directions and
$A_0=\frac c2(u_1u_2^\top+u_2u_1^\top)$, then every quadratic network,
of arbitrary width, with matrix $B$ and intercept $c_0$ satisfies
\[
 \mathbb E[(f_*(X)-q_B(X)-c_0)^2]
 =F_0+2\|B-A_0\|_F^2+c_0^2,
\]
where
\[
 F_0=.64b^2+.08(1-e^{-2})-.64a^4.
\]
The best predictor is realizable by two signed quadratic neurons.
\end{proposition}
\begin{proof}
A centered quadratic form is a linear combination of $Z_i^2-1$ and
$Z_iZ_j$, $i<j$. Independence and oddness imply that $f_*$ has zero
mean and zero inner product with every diagonal term. The product of two
odd hyperbolic tangents has zero correlation with every cross term except
$Z_1Z_2$; for that term the correlation is $.8a^2=c$, while
$\mathbb E[Z_1^2Z_2^2]=1$. The sine term is odd under $Z\mapsto-Z$ and
is orthogonal to every quadratic form and to constants. These facts identify
the projection.

The two components of $f_*$ are orthogonal, with squared norms $.64b^2$
and $.16\mathbb E\sin^2 Z=.08(1-e^{-2})$, respectively; the latter follows
from $\sin^2 z=(1-\cos2z)/2$ and the Gaussian characteristic function.
Subtracting the projection's squared norm $c^2$ yields $F_0$. The
Pythagorean identity and Theorem~\ref{thm:quadratic-geometry} give the risk
formula. Finally, $A_0$ has eigenvectors $(u_1+u_2)/\sqrt2$ and
$(u_1-u_2)/\sqrt2$, eigenvalues $c/2$ and $-c/2$, so two signed neurons
suffice.
\end{proof}
One-dimensional numerical quadrature gives
$a\approx.689027429$, $b\approx.540648381$, $c\approx.379807038$ and
$F_0\approx.111992221$. The added observation-noise variance $.04$ makes
the best total squared risk approximately $.151992221$. These values are
numerical evaluations of the displayed expectations, not additional
statistical estimates. The script \texttt{code/nonlinear\_projection.py}
reproduces the quadrature and evaluates saved final parameters using this
identity. Those post-hoc analytic evaluations are stored separately; all
original Monte Carlo primary outcomes remain unchanged. The sine component alone contributes an irreducible
$.069173177$ for this architecture. More width cannot remove this
misspecification. Information-guided adaptation would need to change the
function family, not merely discover additional quadratic directions.
